\documentclass[11pt, a4paper]{ctexart}

\usepackage{experiment_style}

\begin{document}

% =========================================
% P121. 1
% =========================================
\begin{problem}{P121. 1}
    设
    \[
        A =
        \begin{pmatrix}
            6 & -4 & 18 \\
            20 & -6 & -6 \\
            22 & -22 & 11
        \end{pmatrix},
        \qquad
        \boldsymbol{x}^{(0)}=(1,1,1)^{\mathrm{T}},
    \]
    求 $A$ 的按模最大的特征值和特征向量，要求误差小于 $10^{-6}$。
    \begin{enumerate}[1.]
        \item 直接用幂法计算，观察出现的现象并说明原因；
        \item 用归一化的幂法计算。
    \end{enumerate}
\end{problem}

\begin{solution}
    \textbf{实验内容、步骤及结果}

    幂法函数 \texttt{chap-4/pm.m}：

\begin{lstlisting}[language=Matlab]
function [x_values, lambda_values, iter] = pm(A, x0, max_iter, tol, lambda0, normalize)
    % Power Method for finding the dominant eigenvalue and eigenvector.

    if nargin < 3
        max_iter = inf;
    end

    if nargin < 4
        tol = 0;
    end

    if nargin < 5
        lambda0 = 0;
    end

    if nargin < 6
        normalize = true;
    end

    x = x0;
    lambda = lambda0;
    iter = 0;

    A = A - lambda0 * eye(size(A));

    x_values = x0;
    lambda_values = lambda0;

    while iter < max_iter
        iter = iter + 1;
        if normalize
            [~, idx] = max(abs(x));
            x = x / x(idx);
        end
        x = A * x;

        if normalize
            [~, idx] = max(abs(x));
            lambda_new = x(idx);
        else
            lambda_new = x' * A * x / (x' * x);
        end

        x_values(:, iter + 1) = x / lambda_new;
        lambda_values(iter + 1) = lambda_new + lambda0;

        if tol > 0
            if abs(lambda_new - lambda) < tol
                return;
            end
        end

        lambda = lambda_new;
    end

    if tol > 0
        warning('Power method did not converge within the maximum number of iterations');
    end
end
\end{lstlisting}

    主程序 \texttt{chap-4/P121\_Q1.m}：

\begin{lstlisting}[language=Matlab]
% Define the input matrix A and initial guess x0
A = [6, -4, 18; 20, -6, -6; 22, -22, 11];
x0 = [1; 1; 1];
max_iter = 1e3;
tol = 1e-6;

for normalize = [false, true]
    disp('Normalize:');
    disp(normalize);

    % Call the power method function
    [x_values, lambda_values, iter] = pm(A, x0, max_iter, tol, 0, normalize);

    disp('Number of iterations performed:');
    disp(iter);

    disp('x_values from 100 to 110:');
    disp(x_values(:, 100:110));

    disp('lambda_values from 100 to 110:');
    disp(lambda_values(100:110));

    disp('Latest lambda_value:');
    disp(lambda_values(end));

    disp('----------------------------------------');
end

disp('Eigenvalues of A calculated using MATLAB built-in function:');
disp(eig(A));
\end{lstlisting}

    \textbf{实验结果分析}

\begin{lstlisting}
>> P121_Q1
Normalize:
   0

Warning: Power method did not converge within the maximum number of iterations
> In pm (line 76)
In P121_Q1 (line 12)

Number of iterations performed:
        1000

lambda_values from 100 to 110:
   24.0000   13.8947   24.0000   13.8947   24.0000   13.8947   24.0000   13.8947   24.0000   13.8947   24.0000

Latest lambda_value:
   NaN

----------------------------------------
Normalize:
   1

Warning: Power method did not converge within the maximum number of iterations
> In pm (line 76)
In P121_Q1 (line 12)

Number of iterations performed:
        1000

x_values from 100 to 110:
    1.0000    0.6250    1.0000    0.6250    1.0000    0.6250    1.0000    0.6250    1.0000    0.6250    1.0000
    0.1429    0.6250    0.1429    0.6250    0.1429    0.6250    0.1429    0.6250    0.1429    0.6250    0.1429
    0.5714    1.0000    0.5714    1.0000    0.5714    1.0000    0.5714    1.0000    0.5714    1.0000    0.5714

lambda_values from 100 to 110:
   19.2500   25.1429   19.2500   25.1429   19.2500   25.1429   19.2500   25.1429   19.2500   25.1429   19.2500

Latest lambda_value:
   25.1429

----------------------------------------
Eigenvalues of A calculated using MATLAB built-in function:
  -22.0000
   22.0000
   11.0000
\end{lstlisting}

    直接使用幂法时，迭代向量会随着主导特征值的幂次放大。由于本题矩阵存在按模相等且符号相反的主导特征值 $-22$ 与 $22$，未归一化迭代不仅会发生数值溢出，还会导致特征值估计振荡，最终出现 \texttt{NaN}。

    归一化幂法通过每一步对向量进行缩放避免了溢出，但由于存在两个按模相等的主导特征值，迭代向量和特征值估计仍在两组数值之间交替振荡。因此，本题不能用普通幂法稳定收敛到唯一的主特征向量。可通过位移变换 $A' = A + \mu I$ 打破正负对称主导特征值造成的振荡。
\end{solution}


% =========================================
% P121. 2
% =========================================
\begin{problem}{P121. 2}
    设
    \[
        A =
        \begin{pmatrix}
            12 & 6 & -6 \\
            6 & 16 & 2 \\
            -6 & 2 & 16
        \end{pmatrix},
        \qquad
        x^{(0)}=(1,1,1)^{\mathrm{T}}.
    \]
    先用幂法迭代 3 次，得到 $A$ 的按模最大特征值的近似值，取 $\lambda^\ast$ 为其整数部分，再用反幂法计算 $A$ 的按模最大特征值的更精确近似值，要求误差小于 $10^{-10}$。
\end{problem}

\begin{solution}
    \textbf{实验内容、步骤及结果}

    反幂法函数 \texttt{chap-4/inv\_pm.m}：

\begin{lstlisting}[language=Matlab]
function [x_values, lambda_values, iter] = inv_pm(A, x0, max_iter, tol, lambda0)
    % Inverse Power Method for finding the smallest eigenvalue and eigenvector.

    if nargin < 3
        max_iter = inf;
    end

    if nargin < 4
        tol = 0;
    end

    if nargin < 5
        lambda0 = 0;
    end

    x = x0;
    lambda = lambda0;
    iter = 0;

    A = A - lambda0 * eye(size(A));

    x_values = x0;
    lambda_values = lambda0;

    while iter < max_iter
        iter = iter + 1;
        [~, idx] = max(abs(x));
        x = x / x(idx);
        [L, U] = lu(A);
        y = L \ x;
        x = U \ y;

        [~, idx] = max(abs(x));
        lambda_new = x(idx);

        x_values(:, iter + 1) = x / lambda_new;
        lambda_values(iter + 1) = 1 / lambda_new;

        if tol > 0
            if abs(1 / lambda_new - 1 / lambda) < tol
                return;
            end
        end

        lambda = lambda_new;
    end

    if tol > 0
        warning('Inv power method did not converge within the maximum number of iterations');
    end
end
\end{lstlisting}

    主程序 \texttt{chap-4/P121\_Q2.m}：

\begin{lstlisting}[language=Matlab]
% Define the input matrix A and initial guess x0
A = [12, 6, -6; 6, 16, 2; -6, 2, 16];
x0 = [1; 1; 1];

[x_values_pm, lambda_values_pm, iter_pm] = pm(A, x0, 3);

lambda0 = floor(lambda_values_pm(end));

disp('lambda0:');
disp(lambda0);

[x_values_inv, lambda_values_inv, iter_inv] = inv_pm(A, x0, inf, 1e-10, lambda0);

disp('Number of inv_pm iterations performed:');
disp(iter_inv);

disp('First 3 lambda_values:');
disp(lambda_values_inv(1:3) + lambda0);

disp('Latest x_value:');
disp(x_values_inv(:, end));

disp('Latest lambda_value:');
disp(lambda_values_inv(end) + lambda0);

disp('Eigenvalues of A calculated using MATLAB built-in function:');
disp(eig(A));
\end{lstlisting}

    输出结果如下：

\begin{lstlisting}
>> P121_Q2
lambda0:
    19

Number of inv_pm iterations performed:
    25

First 3 lambda_values:
   38.0000   18.1395   17.8959

Latest x_value:
   -0.0000
    1.0000
    1.0000

Latest lambda_value:
   18.0000

Eigenvalues of A calculated using MATLAB built-in function:
    4.4560
   18.0000
   21.5440
\end{lstlisting}

    使用幂法迭代 3 次后得到位移 $\lambda^\ast = 19$。随后对 $A-\lambda^\ast I$ 使用反幂法，等价于寻找最靠近 $\lambda^\ast$ 的特征值。最终结果收敛到
    \[
        \lambda = 18.0000,
        \qquad
        x \approx (0,1,1)^{\mathrm{T}},
    \]
    与 MATLAB 内置函数 \texttt{eig} 的结果一致。
\end{solution}


% =========================================
% 补充代码
% =========================================
\begin{problem}{补充代码：P120. Q1, Q3, Q4}
\end{problem}

    \texttt{chap-4/P120\_Q1.m}：

\begin{lstlisting}[language=Matlab]
% Define the input matrix A and initial guess x0
A = [-4, 14, 0; -5, 13, 0; -1, 0, 2];
% A = [4, -1, 1; -1, 3, -2; 1, -2, 3];
x0 = [1; 0; 0];
max_iter = 10;

disp('A:');
disp(A);

% Call the power method function
[x_values, lambda_values, iter] = pm(A, x0, max_iter);

disp('Number of iterations performed:');
disp(iter);
disp('First 3 x_values:');
disp(x_values(:, 1:3));
disp('First 3 lambda_values:');
disp(lambda_values(1:3));
disp('Latest x_value:');
disp(x_values(:, end));
disp('Latest lambda_value:');
disp(lambda_values(end));
\end{lstlisting}

    \texttt{chap-4/P120\_Q3.m}：

\begin{lstlisting}[language=Matlab]
% Define the input matrix A and initial guess x0
A = [1, 2, 3; 2, 3, 4; 3, 4, 5];
x0 = [1; 0; 0];
max_iter = inf;
lambda0 = 9.6;

% Call the power method function
[x_values, lambda_values, iter] = inv_pm(A, x0, max_iter, 1e-16, lambda0);

disp('Number of iterations performed:');
disp(iter);
disp('First 3 x_values:');
disp(x_values(:, 1:3));
disp('First 3 lambda_values:');
disp(lambda_values(1:3) + lambda0);
disp('Latest x_value:');
disp(x_values(:, end));
disp('Latest lambda_value:');
disp(lambda_values(end) + lambda0);
\end{lstlisting}

    \texttt{chap-4/P120\_Q4.m}：

\begin{lstlisting}[language=Matlab]
% Define the input matrix A and initial guess x0
A = [2, 1, 0; 1, 3, 1; 0, 1, 4];
x0 = [1; 0; 0];
max_iter = inf;
lambda0 = 1.2679;

% Call the inv power method function
[x_values, lambda_values, iter] = inv_pm(A, x0, max_iter, 1e-16, lambda0);

disp('Number of iterations performed:');
disp(iter);
disp('First 3 x_values:');
disp(x_values(:, 1:3));
disp('First 3 lambda_values:');
disp(lambda_values(1:3) + lambda0);
disp('Latest x_value:');
disp(x_values(:, end));
disp('Latest lambda_value:');
disp(lambda_values(end) + lambda0);
\end{lstlisting}

    运行输出摘要如下：

\begin{lstlisting}
>> P120_Q1
A:
    -4    14     0
    -5    13     0
    -1     0     2

Number of iterations performed:
    10

First 3 x_values:
    1.0000    0.8000    1.0000
         0    1.0000    0.8333
         0    0.2000   -0.0370

First 3 lambda_values:
         0   -5.0000   10.8000

Latest x_value:
    1.0000
    0.7146
   -0.2490

Latest lambda_value:
    6.0084

>> P120_Q1
A:
     4    -1     1
    -1     3    -2
     1    -2     3

Latest x_value:
    1.0000
   -0.9971
    0.9971

Latest lambda_value:
    5.9883

>> P120_Q3
Number of iterations performed:
     8

Latest x_value:
    0.5247
    0.7623
    1.0000

Latest lambda_value:
    9.6235

>> P120_Q4
Number of iterations performed:
     5

Latest x_value:
    1.0000
   -0.7321
    0.2679

Latest lambda_value:
    1.2679
\end{lstlisting}

\end{document}
